\documentclass[11pt]{article}

\usepackage[margin=1in]{geometry}
\usepackage{amsmath,amssymb,bm}
\usepackage{microtype}
\usepackage[hidelinks]{hyperref}

\newcommand{\Tliq}{T_{\mathrm{liq}}}
\newcommand{\evalat}[2]{\left.#1\right|_{#2}}

\title{From Temperature Sensitivities to Melt-Pool Geometry Sensitivities}
\author{}
\date{}

\begin{document}
\maketitle

\section{Liquidus points that define the geometry}

Let $p$ denote one process parameter,
\begin{equation}
  p\in\{P,\sigma,v\},
\end{equation}
and let $T(\bm{x},p)$ be the temperature field for that parameter value.  The
liquidus isotherm is the level set
\begin{equation}
  T(\bm{x},p)=\Tliq.
\end{equation}

The melt-pool width is defined by two geometric extrema of this isotherm:
\begin{itemize}
  \item $y^-(p)$ is the leftmost liquidus point;
  \item $y^+(p)$ is the rightmost liquidus point.
\end{itemize}
Their locations depend on $p$ because changing the process parameter changes
the temperature field and therefore moves the liquidus boundary.  Nevertheless,
both moving points remain on the same liquidus isotherm:
\begin{align}
  T\bigl(\bm{x}^-(p),p\bigr) &= \Tliq, \\
  T\bigl(\bm{x}^+(p),p\bigr) &= \Tliq.
\end{align}
The superscript $\pm$ identifies a spatial extremum; it is not a temperature
maximum or minimum.  Both points have temperature $\Tliq$.

\section{Motion of one liquidus edge}

To expose the essential calculation, suppress the other spatial coordinates
and write the right-edge condition as
\begin{equation}
  T\bigl(y^+(P),P\bigr)=\Tliq.
  \label{eq:right-condition}
\end{equation}
The material is fixed, so $\Tliq$ does not depend on $P$.  Taking the total
derivative of Equation~\eqref{eq:right-condition} gives
\begin{equation}
  \frac{\mathrm{d}}{\mathrm{d}P}
  T\bigl(y^+(P),P\bigr)=0.
\end{equation}
Applying the chain rule,
\begin{equation}
  \evalat{T_y}{y^+}\frac{\partial y^+}{\partial P}
  +\evalat{T_P}{y^+}=0,
\end{equation}
where
\begin{equation}
  T_y=\frac{\partial T}{\partial y},
  \qquad
  T_P=\frac{\partial T}{\partial P}.
\end{equation}
Solving for the motion of the right edge gives
\begin{equation}
  \boxed{
  \frac{\partial y^+}{\partial P}
  =-\evalat{\frac{T_P}{T_y}}{y^+}.}
  \label{eq:right-motion}
\end{equation}

The same calculation at the left edge gives
\begin{equation}
  \boxed{
  \frac{\partial y^-}{\partial P}
  =-\evalat{\frac{T_P}{T_y}}{y^-}.}
  \label{eq:left-motion}
\end{equation}
The signs of $T_y$ will normally differ at the two edges, but
Equations~\eqref{eq:right-motion} and \eqref{eq:left-motion} have the same form.

\section{From edge motion to width sensitivity}

The controller uses the melt-pool \emph{half-width}
\begin{equation}
  w=\frac{y^+-y^-}{2}.
  \label{eq:half-width}
\end{equation}
Differentiating Equation~\eqref{eq:half-width} with respect to $P$ gives
\begin{equation}
  \frac{\partial w}{\partial P}
  =\frac{1}{2}
  \left(
    \frac{\partial y^+}{\partial P}
    -\frac{\partial y^-}{\partial P}
  \right).
\end{equation}
Substituting Equations~\eqref{eq:right-motion} and \eqref{eq:left-motion},
\begin{align}
  \frac{\partial w}{\partial P}
  &=\frac{1}{2}
  \left[
    -\evalat{\frac{T_P}{T_y}}{y^+}
    +\evalat{\frac{T_P}{T_y}}{y^-}
  \right] \\
  &=\boxed{
  \frac{1}{2}
  \left[
    \evalat{\frac{T_P}{T_y}}{y^-}
    -\evalat{\frac{T_P}{T_y}}{y^+}
  \right].}
  \label{eq:width-power}
\end{align}

The factor $1/2$ appears because $w$ is a half-width.  If the reported
quantity were the full width $W=y^+-y^-$, this factor would be absent.

\section{Physical interpretation of the ratio}

At a fixed spatial point,
\begin{itemize}
  \item $T_P$ measures the temperature change produced by a change in power;
  \item $T_y$ measures the temperature change produced by moving in the
        transverse direction.
\end{itemize}
Therefore,
\begin{equation}
  -\frac{T_P}{T_y}
\end{equation}
converts a parameter-induced temperature change into the spatial displacement
required for the point to remain at $T=\Tliq$.  Sampling this ratio at the two
liquidus edges and combining their motions produces the width sensitivity.

\section{General process parameter and depth sensitivity}

Replacing $P$ by any $p\in\{P,\sigma,v\}$ gives
\begin{equation}
  \boxed{
  \frac{\partial w}{\partial p}
  =\frac{1}{2}
  \left[
    \evalat{\frac{T_p}{T_y}}{y^-}
    -\evalat{\frac{T_p}{T_y}}{y^+}
  \right].}
  \label{eq:general-width}
\end{equation}

Let $z^-(p)$ be the deepest liquidus point and define the positive depth by
\begin{equation}
  d=-z^-.
\end{equation}
Implicit differentiation of the liquidus condition at $z^-$ gives
\begin{equation}
  \frac{\partial z^-}{\partial p}
  =-\evalat{\frac{T_p}{T_z}}{z^-}.
\end{equation}
Consequently,
\begin{equation}
  \boxed{
  \frac{\partial d}{\partial p}
  =\evalat{\frac{T_p}{T_z}}{z^-}.}
  \label{eq:general-depth}
\end{equation}

Equations~\eqref{eq:general-width} and \eqref{eq:general-depth}, evaluated for
$p=P,\sigma,v$, form the $2\times3$ geometry Jacobian used by the controller.

\section{Assumptions}

The formulas assume that the relevant liquidus extremum is locally smooth and
unique, and that the required spatial gradient is nonzero: $T_y\neq0$ at the
width edges and $T_z\neq0$ at the deepest point.  In the full three-dimensional
notation, the moving-point condition is
\begin{equation}
  T\bigl(\bm{x}^*(p),p\bigr)=\Tliq,
\end{equation}
and its total derivative is
\begin{equation}
  \nabla T\cdot\frac{\mathrm{d}\bm{x}^*}{\mathrm{d}p}+T_p=0.
\end{equation}
At a smooth support point, the tangency/extremum condition reduces this vector
relation to the one-dimensional edge and depth expressions used above.

\end{document}
